\begin{table}[H]
\centering
\footnotesize
\setlength{\tabcolsep}{4pt}
\begin{tabular}{lrrrrrr}
\toprule
\textbf{$n$} & \textbf{$\rho_S$} & \textbf{$m_S$} & \textbf{$k^*$} & \textbf{$(n-1)\rho_S$} & \textbf{cov at $k$=$n$} & \textbf{se} \\
\midrule
125 & 0.0385 & 7.0 & none & 5 & 0.954 & 0.024 \\
250 & 0.0422 & 14.3 & none & 11 & 0.963 & 0.021 \\
500 & 0.0339 & 21.9 & 398 & 17 & 0.938 & 0.028 \\
1000 & 0.0230 & 32.7 & 980 & 23 & 0.949 & 0.024 \\
2000 & 0.0230 & 40.1 & 1376 & 46 & 0.929 & 0.028 \\
\bottomrule
\end{tabular}
\caption{Sample size does not dissolve the problem. The shipped CausalPFN checkpoint on the synthetic family of Appendix~\ref{app:dgps} at five sample sizes, with subsets drawn uniformly at random, $64$ tasks and $48$ subsets per task. One fixed checkpoint is used throughout, so no training choice of ours varies down the table. $\rho_S$ falls as $n$ grows, but only as roughly $n^{-0.24}$: a sixteenfold increase in $n$ divides it by $1.7$, where the $n^{-1}$ decay that would make the design effect vanish would divide it by $16$. Because the reachable subset size grows with $n$ while $\rho_S$ barely falls, the design effect available at the full sample, $(n-1)\rho_S$, grows from $5$ to $46$ down the table. Coverage at $k=n$ nonetheless stays near nominal, because $m_S$ grows too, at roughly $n^{0.62}$: the checkpoint's per-unit intervals widen relative to its per-unit errors faster than the design effect eats the margin. It is the cushion that holds coverage up here, not the absence of the correlation.}
\label{tab:nsweep}
\end{table}
